索引有三个生命周期阶段：训练生成可复用模型状态，添加编码并排列数据库行，搜索使用完成后的数组。实验必须区分这些阶段。训练状态缓存避免重新学习模型，但仍需要数据库编码与列表构建。冷训练加添加，与复用状态后的查询，测量的是不同操作。

\begin{figure}[tbp]
\centering
% Audited lifecycle; no proposed storage or distributed subsystem.
\begin{tikzpicture}[
 font=\small\linespread{1}\selectfont,
 box/.style={draw=ThesisNavy!65,rounded corners=2pt,align=center,
   text width=3.65cm,minimum height=1.9cm,inner sep=5pt},
 gpu/.style={box,fill=ThesisNavy!7},
 arr/.style={-{Latex[length=2mm]},draw=ThesisNavy!80,thick}]
\node[box] (host) at (0,0) {Host input vectors\\Training samples\\Database batches};
\node[box] (model) at (4.8,0) {Host model\\QR rotation $\Pi$\\Primary codebook};
\node[box] (cache) at (9.6,0) {Host model cache\\$\Pi$ and codebooks\\IVF centroids};
\node[gpu] (train) at (0,-3.6) {GPU training\\Rotate samples\\Fit IVF centroids\\Fit residual model};
\node[gpu] (add) at (4.8,-3.6) {GPU bulk build\\1 or 2 passes\\Assign / encode\\Reorder / pack\\Floats $\leq768$ MiB};
\node[gpu] (index) at (9.6,-3.6) {GPU-resident index\\Offsets and IDs\\Primary / residual codes\\Corrections};
\node[box] (query) at (0,-7.2) {Host query batch\\Pinned staging\\Transfer to GPU};
\node[gpu] (search) at (4.8,-7.2) {GPU query pipeline\\Rotate; probe; scan\\Refine; final top-$k$\\Workspace / graph};
\node[box] (result) at (9.6,-7.2) {Host query results\\IDs and scores\\Transfer from GPU};
\draw[arr] (host) -- node[left,font=\footnotesize] {samples} (train);
\draw[arr] (model.south west) -- (train.north east);
\draw[arr] (train) -- node[above,font=\footnotesize] {model} (add);
\draw[arr] (host.west) -- ++(-.5,0) |- (4.8,-5.15)
 node[pos=.8,below,font=\footnotesize] {database batches} -- (add.south);
\draw[arr] (model) -- (add);
\draw[arr] (add) -- (index);
\draw[<->,dashed,draw=ThesisNavy!80,thick]
 (train.north east) -- node[pos=.8,above,sloped,font=\footnotesize,fill=white,inner sep=1pt]
 {save / load model} (cache.south west);
\draw[arr] (query) -- (search);
\draw[arr] (search) -- (result);
\draw[arr] (index.south) -- ++(0,-.6) -| (search.north);
\end{tikzpicture}

\caption[构建与查询架构]{已实现的建库与查询边界。主机负责采样、正交矩阵构造及解析主码码本生成；评测配置在 GPU 上拟合 IVF 与残差模型。缓存保存训练状态，完成后的数据库表示常驻设备。虚线表示训练状态复用。暂存标注仅约束源码默认的原始／旋转浮点缓冲区，其他流水线 scratch 另计。}
\label{fig:4-architecture}
\end{figure}

主机持有输入向量并提供训练样本，构造方形变换和解析主码码本。GPU 矩阵乘法应用变换；设备训练路径拟合粗粒度 IVF 中心及标量残差码本。添加阶段传输有界批次，编码两级表示，确定 IVF 成员关系，再写成统一的列表顺序表示。因此，主机编排与设备计算共同组成 GPU 常驻搜索系统。

模型状态包含变换、主码码本、残差码本和 IVF 中心。数据库状态进一步包含列表偏移、原始行 ID、主码、残差码及修正项。搜索工作区的逻辑寿命更短：旋转查询、中心分数、探针偏移、LUT 和保留候选缓冲区在下一批被覆盖，但分配的内存保留复用。捕获的 CUDA graph 属于执行状态，引用这些现有数组。

API 通过一次成功的 \texttt{add} 构造固定快照，并拒绝后续 add。搜索该快照时物理位置保持固定。数据变更通过重建替代快照处理，形成简单的批量重建生命周期；可变插入／删除及快照的事务发布还需要额外索引管理。

搜索接收主机查询并返回主机ID及计算分数。分数中的查询常数及其对范围阈值的影响
见\cref{eq:5-refined-constant}，实测主机到主机调用边界见\cref{sec:3-3}。


该设计从主机内存分批传输构建输入，同时将压缩数据库保留在 GPU。它没有提供依赖磁盘的搜索路径，也不意味着任意大的压缩索引都能放入单个设备。
